% संपूर्ण मराठी ग्रंथातील प्रकरण 46: संपूर्णता आणि कॅनॉनिकल प्रतिमाने
% हा संपादनयोग्य स्रोतखंड आहे. संकलनासाठी संपूर्ण स्रोतसंग्रहातील
% अक्षररूपे, पूर्वभाग, चिन्हव्याख्या आणि आकृती-स्रोत आवश्यक आहेत.
% मूळ: Open Logic Project; मजकूर आणि रूपांतर CC BY 4.0.
% यंत्रानुवाद, दुरुस्ती आणि तपासणी: OpenAI Codex —
% GPT-5.6 Sol आणि GPT-6 Sol, दोन्ही Ultra effort.
% दुरुस्ती-पुनरावलोकन: GPT-6.1 Sol, Ultra effort.
\chapter{संपूर्णता आणि कॅनॉनिकल प्रतिमाने}\label{nml:com:chap}










\section{प्रस्तावना}\label{nml:com:int:sec}

$\Sigma$ ही मोडल प्रणाली असेल, तर निर्दोषता
प्रमेय असे सांगते: $\Sigma \Proves \varphi$ असल्यास,
ज्या प्रतिमान-वर्ग~$\mClass{C}$ मध्ये~$\Sigma$
मधील सर्व सूत्रांचे सर्व दृष्टांत वैध
आहेत, त्या वर्गात $\varphi$ वैध असते.
विशेषतः, $\Log{K}
\Proves \varphi$ असेल, तर $\varphi$ सर्व प्रतिमानांत
सत्य असते; $\Log{KT} \Proves \varphi$ असेल,
तर $\varphi$ सर्व परावर्ती प्रतिमानांत सत्य असते;
$\Log{KD} \Proves \varphi$ असेल, तर $\varphi$
सर्व सर्वत्र उत्तरवर्ती प्रतिमानांत सत्य असते; इत्यादी.

संपूर्णता हा निर्दोषतेचा व्यत्यास आहे:
उदाहरणार्थ, \Log{K} संपूर्ण आहे याचा अर्थ
सूत्र~$\varphi$ वैध असेल, तर $\Proves \varphi$.
संपूर्णता सिद्ध करणे निर्दोषता सिद्ध करण्यापेक्षा
खूप कठीण आहे. सुरुवातीला प्रतिपरिवर्तन विचारात
घेणे उपयुक्त ठरते: $\Proves/ \varphi$ असेल तेव्हा
नेहमी प्रतिप्रतिमान, म्हणजे
$\mSat/{M}{\varphi}$ असे~$\mModel{M}$ प्रतिमान,
अस्तित्वात असेल, तर आणि तेव्हाच \Log{K}
संपूर्ण असते. समतुल्यपणे ($\varphi$ चा निषेध
घेतल्यास), $\Proves/ \lnot \varphi$ असेल तेव्हा
$\varphi$ चे प्रतिमान असते, हे सिद्ध करता येईल.
असे प्रतिमान रचताना $\varphi$ मधील माहिती
वापरता येते. विशिष्ट सूत्रे साठी
प्रतिमाने शोधतानाही आपण अनेकदा असेच करतो.
उदाहरणार्थ, $p \lif \Box q$ चे प्रतिप्रतिमान
शोधायचे असेल, तर त्यात $p$ सत्य आणि
$\Box q$ असत्य असलेले जग हवे, हे माहीत असते.
ज्या जगात $\Box q$ असत्य आहे, त्या जगापासून
प्राप्य असे एखादे जग असले पाहिजे की त्यात $q$
असत्य आहे. आपल्याला एवढेच ठरवायचे असते:
कोणत्या जगांत विधानीय चले
सत्य आहेत आणि कोणती जगे कोणत्या जगांपासून
प्राप्य आहेत.

पण संपूर्णता सिद्ध करताना प्रतिमान रचण्यासाठी
आपल्यापाशी एखादे ठरावीक सूत्र~$\varphi$
नसते. $\Proves/[\Sigma] \lnot \varphi$ असलेल्या
प्रत्येक $\varphi$ साठी प्रतिमान अस्तित्वात आहे,
हे दाखवायचे असते. ही किमान आवश्यक अट आहे:
\emph{जर} $\Proves[\Sigma] \lnot \varphi$ असेल,
तर निर्दोषतेनुसार~$\varphi$ चे
($\Sigma$ सत्य असलेले) प्रतिमान असू शकत नाही.
आता लक्षात घ्या की $\Proves/[\Sigma]
\lnot \varphi$ असेल तेव्हा आणि तेव्हाच $\varphi$ हे
$\Sigma$-सुसंगत असते. ($\Sigma
\Proves/[\Sigma] \lnot \varphi$ आणि
$\varphi \Proves/[\Sigma] \lfalse$ समतुल्य
आहेत, हे आठवा.) म्हणून प्रत्येक
$\Sigma$-सुसंगत सूत्र साठी प्रतिमान
रचणे हे आपले काम आहे.

आपण पुढील युक्ती वापरू: $\varphi$ सामावणारा
$\Sigma$-सुसंगत सूत्रांचा संच शोधू;
त्यात $\varphi$ सत्य करणारे जग कसे असावे
हे सांगणारी इतर सूत्रेही असतील. असे संच
\emph{संपूर्ण} $\Sigma$-सुसंगत संच असतात.
प्रत्येक सुसंगत $\varphi$ साठी काम करणाऱ्या सर्वसाधारण रचनेत
फक्त एकच जग ठेवणे नेहमी पुरेसे नसते. अनेक जगे आणि
त्यांच्यातील प्राप्यता संबंध असण्यास रचनेने परवानगी दिली पाहिजे. $\varphi$ सामावणाऱ्या
संपूर्ण $\Sigma$-सुसंगत संचात
$\Box \psi$ आणि~$\Diamond \chi$ या रूपांतील
इतर सूत्रे ही असतील. सर्व प्राप्य
जगांत $\psi$ सत्य असले पाहिजे; किमान एका
प्राप्य जगात $\chi$ सत्य असले पाहिजे. यासाठी
शक्य असलेले \emph{सर्व} संपूर्ण
$\Sigma$-सुसंगत संच जगांच्या संचाचा
आधार म्हणून घेऊ. आपल्या प्रतिमानात एक
संपूर्ण $\Sigma$-सुसंगत संच दुसऱ्यापासून
केव्हा प्राप्य मानायचा, हे ठरवणे अवघड
ठरेल.

असे परिभाषित केलेल्या प्रतिमानात एखाद्या
जगात---जे स्वतः संपूर्ण $\Sigma$-सुसंगत
संच आहे---$\varphi$ सत्य असेल तेव्हा आणि
तेव्हाच $\varphi$ त्या संचाचा घटक
असेल, हे दाखवू. $\varphi$ हे
$\Sigma$-सुसंगत असेल, तर किमान एका
संपूर्ण $\Sigma$-सुसंगत संचाचा ते
घटक असेल (हे आपण सिद्ध करू);
म्हणून $\varphi$~सत्य असलेले जग मिळेल. अशा
रीतीने प्रत्येक $\Sigma$-सुसंगत
सूत्र~$\varphi$ काहीतरी जगात सत्य असलेले
एकच प्रतिमान आपल्याला मिळेल.
हेच~$\Sigma$ साठीचे \emph{कॅनॉनिकल}
प्रतिमान आहे.













\section{संपूर्ण $\Sigma$-सुसंगत संच}\label{nml:com:ccs:sec}

$\Sigma$ हा मोडल सूत्रांचा संच आहे असे
समजा---सामान्य मोडल तर्कशास्त्राची स्वयंसिद्धके
किंवा परिभाषक तत्त्वे म्हणून त्यांच्याकडे पाहा.
संच~$\Gamma$ हा $\Sigma$-सुसंगत असेल तेव्हा
आणि तेव्हाच $\Gamma \Proves/[\Sigma] \lfalse$;
म्हणजे प्रत्येक $\varphi_i \in \Gamma$ असताना
$\Sigma$ पासून~$\varphi_1 \lif (\varphi_2 \lif \cdots (\varphi_n \lif
\lfalse)\dots)$ ची निष्पत्ती नसते.
आपण ``कॅनॉनिकल'' प्रतिमान रचू; त्यात प्रत्येक
जग विशिष्ट प्रकारचा $\Sigma$-सुसंगत संच
असेल. तो केवळ $\Sigma$-सुसंगत नसून
महत्तम अर्थाने तसा असेल: प्रत्येक मोडल
सूत्राचे सत्यत्व तो ठरवेल.
प्रत्येक $\varphi$ साठी $\varphi
\in \Gamma$ किंवा $\lnot \varphi \in \Gamma$:

\begin{defn}
  संच~$\Gamma$ हा \emph{संपूर्ण $\Sigma$-सुसंगत}
  असेल तेव्हा आणि तेव्हाच तो $\Sigma$-सुसंगत
  असेल आणि प्रत्येक~$\varphi$ साठी
  $\varphi \in \Gamma$ किंवा $\lnot \varphi \in \Gamma$
  यांपैकी एक विधान खरे असेल.
\end{defn}

संपूर्ण $\Sigma$-सुसंगत संच~$\Gamma$ चे अनेक
उपयुक्त गुणधर्म आहेत. प्रथम, ते निगामीदृष्ट्या
बंद असतात: $\Gamma
\Proves[\Sigma] \varphi$ असेल, तर $\varphi \in \Gamma$.
विशेषतः, सूत्र~$\varphi \in \Sigma$ चा प्रत्येक
दृष्टांतही $\in
\Gamma$ असतो. शिवाय $\Gamma$ मधील सदस्यत्व
विधानीय संयोगकांच्या सत्यत्व-अटींशी जुळते.
आपण ``कॅनॉनिकल प्रतिमान'' परिभाषित करू तेव्हा
हे महत्त्वाचे ठरेल.

\begin{prop}\label{nml:com:ccs:prop:ccs-properties}
  $\Gamma$ हा संपूर्ण $\Sigma$-सुसंगत आहे असे
  समजा. मग:
  \begin{enumerate}
  \item \label{nml:com:ccs:prop:ccs-closed}
    $\Gamma$ हा~$\Sigma$ च्या सापेक्ष
    निगामीदृष्ट्या बंद आहे.
  \item \label{nml:com:ccs:prop:ccs-sigma}
    $\Sigma \subseteq \Gamma$.
  \item \label{nml:com:ccs:prop:ccs-lfalse}
    $\lfalse \notin \Gamma$
  \item \label{nml:com:ccs:prop:ccs-ltrue}
    $\ltrue \in \Gamma$
  \item \label{nml:com:ccs:prop:ccs-lnot}
    $\lnot\varphi \in \Gamma$ असेल तेव्हा आणि
    तेव्हाच $\varphi \notin
    \Gamma$.
  \item \label{nml:com:ccs:prop:ccs-land}
    $\varphi \land \psi \in \Gamma$ असेल तेव्हा आणि तेव्हाच
    $\varphi \in \Gamma$ आणि $\psi \in \Gamma$ असतील.
  \item \label{nml:com:ccs:prop:ccs-lor}
    $\varphi \lor \psi \in \Gamma$ असेल तेव्हा आणि तेव्हाच
    $\varphi \in \Gamma$ किंवा $\psi \in \Gamma$ असेल.
  \item \label{nml:com:ccs:prop:ccs-lif}
    $\varphi \lif \psi \in \Gamma$ असेल तेव्हा आणि तेव्हाच
    $\varphi \notin \Gamma$ किंवा $\psi \in \Gamma$ असेल.
  \item \label{nml:com:ccs:prop:ccs-liff}
    $\varphi \liff \psi \in \Gamma$ असेल तेव्हा आणि तेव्हाच
    $\varphi \in \Gamma$ आणि $\psi \in
    \Gamma$ दोन्ही असतील, किंवा
    $\varphi \notin \Gamma$ आणि $\psi \notin \Gamma$
    दोन्ही असतील.
  \end{enumerate}
\end{prop}

\begin{proof}
  \begin{enumerate}
  \item $\Gamma \Proves[\Sigma] \varphi$ पण $\varphi \notin
    \Gamma$ असे समजा. $\Gamma$ संपूर्ण
    $\Sigma$-सुसंगत असल्याने
    $\lnot\varphi \in \Gamma$. मग $\Gamma$ विसंगत
    ठरेल, कारण $\varphi, \lnot \varphi \Proves[\Sigma] \lfalse$.

  \item $\varphi \in \Sigma$ असेल, तर
    $\Gamma \Proves[\Sigma] \varphi$; आणि
    निगामी बंदतेवरून, म्हणजे
    \ref{nml:com:ccs:prop:ccs-closed} वरून, $\varphi
    \in \Gamma$.

  \item $\lfalse \in \Gamma$ असेल,
    तर $\Gamma \Proves[\Sigma] \lfalse$;
    त्यामुळे $\Gamma$ हा $\Sigma$-विसंगत
    ठरेल.

  \item $\Gamma \Proves[\Sigma] \ltrue$;
    म्हणून निगामी बंदतेवरून, म्हणजे
    \ref{nml:com:ccs:prop:ccs-closed} वरून,
    $\ltrue \in \Gamma$.

\item $\lnot\varphi \in \Gamma$ असेल, तर
    सुसंगततेमुळे $\varphi \notin \Gamma$.
    उलट $\varphi \notin \Gamma$ असेल, तर
    $\Gamma$ संपूर्ण $\Sigma$-सुसंगत
    असल्याने $\lnot\varphi \in \Gamma$.

  \item $\varphi \land \psi \in
      \Gamma$ असे समजा. $(\varphi \land \psi) \lif \varphi$
      हा सर्वतः-सत्य दृष्टांत असल्याने,
      निगामी बंदतेवरून, म्हणजे
      \ref{nml:com:ccs:prop:ccs-closed} वरून, $\varphi \in \Gamma$.
      $\psi \in \Gamma$ साठीही तसेच. उलट
      $\varphi \in \Gamma$ आणि $\psi \in
      \Gamma$ दोन्ही असतील, तर निगामी
      बंदतेवरून $(\varphi \land \psi) \in
      \Gamma$; कारण $\varphi \lif (\psi \lif (\varphi \land \psi))$
      हा सर्वतः-सत्य दृष्टांत आहे.

  \item $\varphi \lor \psi \in
      \Gamma$, पण $\varphi \notin \Gamma$ आणि
      $\psi \notin \Gamma$ असे समजा. $\Gamma$
      संपूर्ण $\Sigma$-सुसंगत असल्याने
      $\lnot\varphi \in \Gamma$ आणि
      $\lnot \psi \in \Gamma$. मग
      $\lnot(\varphi \lor \psi) \in \Gamma$;
      कारण $\lnot \varphi \lif (\lnot \psi \lif \lnot (\varphi \lor \psi))$
      हा सर्वतः-सत्य दृष्टांत आहे. त्यामुळे
      $\Gamma$ हा $\Sigma$-विसंगत ठरेल,
      हा व्याघात आहे. उलट $\varphi \in \Gamma$ असेल,
      तर $\varphi \lif (\varphi \lor \psi)$ हा सर्वतः-सत्य दृष्टांत
      आणि निगामी बंदता यांवरून $\varphi \lor \psi \in \Gamma$.
      तसेच $\psi \in \Gamma$ असेल, तर
      $\psi \lif (\varphi \lor \psi)$ वापरून तोच निष्कर्ष मिळतो.

  \item $\varphi \lif \psi \in
      \Gamma$ आणि $\varphi \in \Gamma$ असे समजा.
      मग $\Gamma \Proves[\Sigma] \psi$;
      त्यामुळे निगामी बंदतेवरून $\psi \in \Gamma$.
      उलट $\varphi
      \lif \psi \notin \Gamma$ असेल, तर
      $\Gamma$ संपूर्ण $\Sigma$-सुसंगत
      असल्याने $\lnot (\varphi \lif \psi) \in \Gamma$.
      $\lnot(\varphi \lif \psi) \lif \varphi$ हा
      सर्वतः-सत्य दृष्टांत असल्याने निगामी
      बंदतेवरून $\varphi \in
      \Gamma$. तसेच $\lnot(\varphi \lif \psi) \lif
      \lnot \psi$ हा सर्वतः-सत्य दृष्टांत असल्याने
      $\lnot \psi \in
      \Gamma$. $\Gamma$ हा $\Sigma$-सुसंगत
      असल्यामुळे $\psi \notin \Gamma$.

  \item $\varphi \liff \psi \in
      \Gamma$ असे समजा. $\varphi \in \Gamma$
      असेल, तर $\psi \in \Gamma$; कारण $(\varphi
      \liff \psi) \lif (\varphi \lif \psi)$ हा
      सर्वतः-सत्य दृष्टांत आहे. तसेच
      $\psi \in \Gamma$ असेल, तर $\varphi \in \Gamma$.
      म्हणून $\varphi \in \Gamma$ आणि $\psi \in \Gamma$
      दोन्ही असतील, किंवा $\varphi \in \Gamma$
      आणि $\psi \in \Gamma$ यांपैकी कोणतेही नसेल.

      उलट $\varphi \liff \psi \notin \Gamma$
      असे समजा. $\Gamma$ संपूर्ण
      $\Sigma$-सुसंगत असल्याने
      $\lnot (\varphi \liff \psi) \in
      \Gamma$. $\lnot(\varphi \liff \psi) \lif (\varphi \lif \lnot \psi)$
      हा सर्वतः-सत्य दृष्टांत असल्याने
      $\varphi \in \Gamma$ असेल, तर
      $\lnot \psi \in
      \Gamma$; आणि $\Gamma$ हा
      $\Sigma$-सुसंगत असल्याने $\psi \notin
      \Gamma$. तसेच $\psi \in \Gamma$
      असेल, तर $\varphi \notin \Gamma$.
      म्हणून $\varphi \in \Gamma$ आणि $\psi \in \Gamma$
      ही दोन्ही विधाने एकत्र खरी असण्याची शक्यता नाही;
      तसेच $\varphi \notin \Gamma$ आणि $\psi \notin \Gamma$
      ही दोन्ही विधाने एकत्र खरी असण्याचीही शक्यता नाही.
  \end{enumerate}
\end{proof}















\section{लिंडेनबाउमचे पूर्वप्रमेय}\label{nml:com:lin:sec}

लिंडेनबाउमचे पूर्वप्रमेय असे सांगते की सूत्रांचा प्रत्येक
$\Sigma$-सुसंगत संच एखाद्या \emph{संपूर्ण} $\Sigma$-सुसंगत
संचात समाविष्ट असतो. कॅनॉनिकल प्रतिमानाची आपली रचना दाखवेल की
प्रत्येक संपूर्ण $\Sigma$-सुसंगत संच~$\Delta$ साठी कॅनॉनिकल
प्रतिमानात असे एक जग असते की नेमके~$\Delta$ मधील
सूत्रे त्या जगात सत्य असतात. म्हणून लिंडेनबाउमचे पूर्वप्रमेय
प्रत्येक $\Sigma$-सुसंगत संच कॅनॉनिकल प्रतिमानातील कोणत्या तरी
जगात सत्य असल्याची हमी देते.

\begin{thm}[लिंडेनबाउमचे पूर्वप्रमेय]\label{nml:com:lin:thm:lindenbaum}
  $\Gamma$ हा $\Sigma$-सुसंगत असेल, तर~$\Gamma$ चा विस्तार असलेला
  संपूर्ण $\Sigma$-सुसंगत संच~$\Delta$ अस्तित्वात असतो.
\end{thm}

\begin{proof}
भाषेतील सर्व सूत्रांची $\varphi_0$, $\varphi_1$, \dots{} अशी यादी असू द्या;
त्यात एखादे सूत्र पुन्हा येऊ शकते. उदाहरणार्थ, आधी
$\Obj p_0$ नोंदवा आणि नंतर प्रत्येक $n \ge 1$ टप्प्यावर
$\Obj p_0$, \dots,~$\Obj p_n$ यांमधीलच चल वापरून
बनणारी व लांबी~$n$ पर्यंतची सर्व सूत्रे नोंदवा;
अशा सूत्रांची संख्या सांत असते. आपण~$n$ वरील विगमनाने
सूत्रांचे संच~$\Delta_n$ परिभाषित करू आणि मग
$\Delta = \bigcup_n \Delta_n$ ठेवू. प्रथम $\Delta_0 = \Gamma$
ठेवू. $\Delta_n$ परिभाषित झाला आहे असे धरून
$\Delta_{n+1}$ ची व्याख्या अशी करू:
\[\Delta_{n+1} =
\begin{cases}
  \Delta_n \cup \{\varphi_n\}, & \text{जर $\Delta_n \cup \{ \varphi_n\}$
    हा $\Sigma$-सुसंगत असेल;} \\
  \Delta_n \cup \{ \lnot \varphi_n\}, & \text{अन्यथा.}
\end{cases}
\]
आता $\Delta = \bigcup_{n=0}^\infty \Delta_n$ ठेवू.

या व्याख्येतून आवश्यक गुणधर्म असलेला संच~$\Delta$
खरोखर मिळतो हे दाखवायचे आहे: $\Gamma \subseteq \Delta$
आणि~$\Delta$ हा संपूर्ण $\Sigma$-सुसंगत आहे.

$\Gamma \subseteq \Delta$ हे उघड आहे, कारण रचनेनुसार
$\Delta_0 \subseteq \Delta$ आणि $\Delta_0 = \Gamma$.
खरे तर प्रत्येक~$n$ साठी $\Delta_n \subseteq \Delta$,
कारण~$\Delta$ हा सर्व $\Delta_n$ चा संघ आहे. (रचनेच्या
प्रत्येक पायरीवर आधीच्या संचात सूत्र जोडतो; म्हणून
$\Delta_n \subseteq \Delta_{n+1}$. तसेच
$\subseteq$ संक्रमणशील असल्याने $n \le m$ असेल तेव्हा
$\Delta_n \subseteq \Delta_{m}$.) रचनेच्या प्रत्येक
टप्प्यावर आपण $\varphi_n$ किंवा $\lnot \varphi_n$ जोडतो आणि
प्रत्येक सूत्र हे सर्व~$\varphi_n$ च्या यादीत किमान
एकदा येते. म्हणून प्रत्येक $\varphi$ साठी $\varphi
\in \Delta$ किंवा $\lnot \varphi \in \Delta$;
त्यामुळे व्याख्येनुसार~$\Delta$ संपूर्ण आहे.

शेवटी~$\Delta$ हा $\Sigma$-सुसंगत आहे हे दाखवायचे आहे.
त्यासाठी (a)~$\Delta$ हा $\Sigma$-विसंगत असता, तर
एखादा~$\Delta_n$ $\Sigma$-विसंगत असता, आणि
(b)~सर्व $\Delta_n$ हे $\Sigma$-सुसंगत आहेत,
असे आपण दाखवू.

समजा~$\Delta$ हा $\Sigma$-विसंगत आहे. मग
$\Delta \Proves[\Sigma] \lfalse$; म्हणजे
$\varphi_1$, \dots,~$\varphi_k \in \Delta$ अशी सूत्रे असतात की
$\Sigma \Proves \varphi_1 \lif (\varphi_2 \lif \cdots (\varphi_k
\lif \lfalse)\dots)$. $\Delta = \bigcup_{n=0}^\infty \Delta_n$
असल्याने प्रत्येक $\varphi_i \in \Delta_{n_i}$ असे कोणत्यातरी~$n_i$
साठी असते. $n=\max(\{0\}\cup\{n_1,\dots,n_k\})$ घ्या;
यामुळे शून्य गृहीतकांच्या $k=0$ प्रकरणातही $n=0$ ठरतो. $n_i \le n$ असल्यामुळे
$\Delta_{n_i} \subseteq \Delta_n$. म्हणून ही सर्व
$\varphi_i$ सूत्रे एखाद्या~$\Delta_n$ मध्ये असतात.
याचा अर्थ $\Delta_n \Proves[\Sigma] \lfalse$,
म्हणजेच~$\Delta_n$ हा $\Sigma$-विसंगत आहे.

प्रत्येक $\Delta_n$ $\Sigma$-सुसंगत आहे हे दाखवण्यासाठी
आपण~$n$ वरील सोपे विगमन वापरू. $\Delta_0 = \Gamma$
आणि~$\Gamma$ $\Sigma$-सुसंगत असल्याचे गृहीत आहे;
म्हणून $n = 0$ साठी दावा खरा आहे. आता तो~$n$
साठी खरा आहे, म्हणजे $\Delta_n$ $\Sigma$-सुसंगत
आहे, असे समजा. $\Delta_n \cup \{\varphi_n\}$ हा
$\Sigma$-सुसंगत असेल, तर $\Delta_{n+1}$ हाच
संच असतो; अन्यथा तो $\Delta_n \cup \{\lnot\varphi_n\}$
असतो. पहिल्या प्रकरणात $\Delta_{n+1}$
$\Sigma$-सुसंगत असणे स्पष्ट आहे. परंतु
\ref{nml:prf:con:prop:consistencyfacts}\ref{nml:prf:con:prop:consistencyfacts-c}
नुसार $\Delta_n \cup \{\varphi_n\}$ किंवा
$\Delta_n \cup \{\lnot\varphi_n\}$ यांपैकी एक
सुसंगत असतो; म्हणून दुसऱ्या प्रकरणातही
$\Delta_{n+1}$ सुसंगत असतो.
\end{proof}

\begin{cor}\label{nml:com:lin:cor:provability-characterization}
  $\Gamma \Proves[\Sigma] \varphi$ असेल तेव्हा आणि तेव्हाच
  $\Gamma$ चा विस्तार असलेल्या प्रत्येक संपूर्ण
  $\Sigma$-सुसंगत संच~$\Delta$ मध्ये $\varphi \in \Delta$
  असते. ($\Gamma = \emptyset$ असतानाही हे लागू होते;
  त्या बाबतीत मोडल प्रणाली~$\Sigma$ चे आणखी एक
  लक्षणन मिळते.)
\end{cor}

\begin{proof}
  $\Gamma \Proves[\Sigma] \varphi$ असे समजा आणि
  $\Gamma$ चा विस्तार असलेला कोणताही संपूर्ण
  $\Sigma$-सुसंगत संच~$\Delta$ घ्या. जर $\varphi
  \notin \Delta$ असेल, तर महत्तमतेमुळे
  $\lnot\varphi \in \Delta$. त्यामुळे एकस्वनिकतेने
  $\Delta \Proves[\Sigma] \varphi$ आणि स्वतःच्या गृहीतकावरून
  $\Delta \Proves[\Sigma] \lnot\varphi$; म्हणून~$\Delta$
  विसंगत ठरेल. उलट $\Gamma \Proves/[\Sigma] \varphi$
  असेल, तर $\Gamma \cup \{ \lnot\varphi\}$ हा
  $\Sigma$-सुसंगत संच असतो. लिंडेनबाउमच्या
  पूर्वप्रमेयानुसार $\Gamma \cup \{ \lnot\varphi \}$
  चा विस्तार असलेला संपूर्ण सुसंगत संच~$\Delta$
  अस्तित्वात असतो. सुसंगततेमुळे $\varphi
  \notin \Delta$.
\end{proof}













\section{मोडल कारक आणि संपूर्ण सुसंगत संच}\label{nml:com:mod:sec}

\begin{explain}
एखाद्या सामान्य मोडल तर्कशास्त्र~$\Sigma$ मधील संपूर्ण
$\Sigma$-सुसंगत संच~$\Delta$ हे ज्याचे जग आहेत,
असे प्रतिमान~$\mModel{M^\Sigma}$ रचताना अशा “जगां”मधील
प्राप्यता संबंध~$R^\Sigma$ देखील परिभाषित करावा लागेल.
प्राप्यता संबंध (आणि मूल्यांकन~$V^\Sigma$) यांची व्याख्या
अशी हवी की $\mSat{M^\Sigma}{\varphi}[\Delta]$ असेल तेव्हा
आणि तेव्हाच $\varphi \in \Delta$. हे कसे करावे?

प्राप्यता संबंध परिभाषित केल्यावर, जगातील सत्यतेच्या
व्याख्येनुसार $\mSat{M^\Sigma}{\Box\varphi}[\Delta]$
  असेल तेव्हा आणि तेव्हाच $R^\Sigma\Delta\Delta'$ असलेल्या
  प्रत्येक $\Delta'$ साठी $\mSat{M^\Sigma}{\varphi}[\Delta']$.
$\mSat{M^\Sigma}{\varphi}[\Delta]$ असेल तेव्हा आणि तेव्हाच
$\varphi \in \Delta$ हे सिद्ध करण्यासाठी हे विशेषतः
मोडल कारकाने सुरू होणाऱ्या सूत्रे साठी खरे असणे
आवश्यक आहे; म्हणजे
$\mSat{M^\Sigma}{\Box\varphi}[\Delta]$ असेल
  तेव्हा आणि तेव्हाच $\Box \varphi \in \Delta$.
ही आवश्यकता आणि
$\Box \varphi$
साठी जगातील सत्यतेची व्याख्या एकत्र केल्यास:
\[  \Box\varphi \in \Delta \text{ तेव्हा आणि तेव्हाच } \varphi \in \Delta' \text{
    प्रत्येक $\Delta'$ साठी, जेथे } R^\Sigma\Delta\Delta'
\]
डावीकडून उजवीकडे जाणारी दिशा पाहा:
  $\Box \varphi \in \Delta$ असेल, तर कोणत्याही $\varphi$
  आणि $R^\Sigma\Delta\Delta'$ असलेल्या कोणत्याही
  $\Delta'$ साठी $\varphi \in \Delta'$ असे ती सांगते.
  प्रत्येक $\Box\varphi \in \Delta$ साठी $\varphi \in \Delta'$
  असेल तेव्हा आणि तेव्हाच $R^\Sigma\Delta\Delta'$
  अशी अट ठेवल्यास हे खरे ठरते. या द्विसशर्ताच्या
  उजव्या बाजूची अट संक्षिप्तपणे
  $\Setabs{\varphi}{\Box\varphi \in \Delta} \subseteq \Delta'$
  अशी लिहिता येते.

मग प्रश्न असा आहे: $R^\Sigma$ ची ही व्याख्या
$\Box \varphi \in \Delta$ असेल तेव्हा आणि
  तेव्हाच $\mSat{M^\Sigma}{\Box \varphi}[\Delta]$ याची
  खरोखर हमी देते का? ती
    $\Diamond \varphi \in \Delta$
  असेल तेव्हा आणि तेव्हाच
  $\mSat{M^\Sigma}{\Diamond \varphi}[\Delta]$ याचीही
  हमी देते का? पुढील काही निष्कर्ष हे सिद्ध करतील.
\end{explain}

\begin{defn}
  $\Gamma$ हा सूत्रांचा संच असेल, तर
  \begin{align*}
    \Box\Gamma & = \Setabs{\Box \psi}{\psi \in \Gamma}\\
    \Diamond\Gamma & = \Setabs{\Diamond \psi}{\psi \in \Gamma}\\
    \intertext{आणि}
    \Box^{-1}\Gamma & = \Setabs{\psi}{\Box \psi \in \Gamma}\\
    \Diamond^{-1}\Gamma & = \Setabs{\psi}{\Diamond \psi \in \Gamma}\\
  \end{align*}
\end{defn}

दुसऱ्या शब्दांत, $\Box\Gamma$ म्हणजे~$\Gamma$ मधील
प्रत्येक सूत्राच्या पुढे $\Box$ लावलेला संच;
ही सर्व सूत्रे~$\Gamma$ मधूनच घेतली जातात.
$\Box^{-1}\Gamma$ मध्ये~$\Gamma$ मधील $\Box$ ने
सुरू होणाऱ्या सर्व सूत्रे मधील सुरुवातीचा
$\Box$ काढून टाकला जातो. ही व्याख्या स्वतंत्रपणे
फारशी महत्त्वाची नाही, पण संकेतलेखन बरेच सोपे करते.

लक्षात घ्या की $\Box\Box^{-1}\Gamma \subseteq \Gamma$:
\[\Box\Box^{-1}\Gamma = \Setabs{\Box\psi}{\Box\psi \in \Gamma}
\]
म्हणजे हा~$\Gamma$ मधील $\Box$ ने सुरू होणाऱ्या
सर्व सूत्रांचाच संच आहे.

\begin{lem}\label{nml:com:mod:lem:box1}
  $\Gamma \Proves[\Sigma] \varphi$ असेल, तर
  $\Box\Gamma \Proves[\Sigma] \Box \varphi$.
\end{lem}

\begin{proof}
  $\Gamma \Proves[\Sigma] \varphi$ असेल, तर
  $\psi_1$, \dots, $\psi_k \in \Gamma$ अशी सूत्रे
  असतात. $k=0$ असेल, तर $\Sigma \Proves \varphi$;
  \Nec{} ने $\Sigma \Proves \Box\varphi$ मिळते आणि
  शून्य गृहीतकांच्या प्रकरणानुसार $\Box\Gamma \Proves[\Sigma] \Box\varphi$.
  आता $k\geq1$ असे समजा. मग $\Sigma \Proves \psi_1 \lif (\psi_2 \lif \cdots
  (\psi_k \lif \varphi)\cdots)$. $\Sigma$ सामान्य असल्याने
  \RK{} नियमानुसार $\Sigma \Proves \Box\psi_1 \lif
  (\Box\psi_2 \lif \cdots (\Box\psi_k \lif \Box\varphi)\cdots)$;
  येथे अर्थातच $\Box\psi_1$, \dots, $\Box\psi_k \in
  \Box\Gamma$. म्हणून व्याख्येनुसार
  $\Box\Gamma \Proves[\Sigma] \Box \varphi$.
\end{proof}

\begin{lem}\label{nml:com:mod:lem:box2}
  $\Box^{-1} \Gamma \Proves[\Sigma] \varphi$ असेल,
  तर $\Gamma \Proves[\Sigma] \Box\varphi$.
\end{lem}

\begin{proof}
  $\Box^{-1}\Gamma \Proves[\Sigma] \varphi$ असे
  समजा. मग \ref{nml:com:mod:lem:box1} नुसार
  $\Box\Box^{-1}\Gamma \Proves[\Sigma] \Box\varphi$.
  पण $\Box\Box^{-1}\Gamma \subseteq \Gamma$
  असल्याने एकस्वनिकतेनुसार
  $\Gamma \Proves[\Sigma] \Box\varphi$.
\end{proof}




\begin{prop}\label{nml:com:mod:prop:box}
  $\Gamma$ हा संपूर्ण $\Sigma$-सुसंगत असेल,
  तर $\Box \varphi \in \Gamma$ असेल तेव्हा आणि
  तेव्हाच $\Box^{-1} \Gamma \subseteq \Delta$
  असलेल्या प्रत्येक संपूर्ण $\Sigma$-सुसंगत
  $\Delta$ साठी $\varphi \in \Delta$ असेल.
\end{prop}

\begin{proof}
  $\Gamma$ हा संपूर्ण $\Sigma$-सुसंगत आहे
  असे समजा. डावीकडून उजवीकडे जाणारी दिशा सोपी आहे:
  $\Box \varphi \in \Gamma$ आणि
  $\Box^{-1}\Gamma \subseteq \Delta$ असे
  समजा. $\Box\varphi \in \Gamma$ असल्याने
  $\varphi \in \Box^{-1}\Gamma \subseteq \Delta$;
  म्हणून $\varphi \in \Delta$.

  उजवीकडून डावीकडे जाणाऱ्या दिशेसाठी आपण प्रतिपरिवर्तन सिद्ध करू:
  $\Box\varphi \notin \Gamma$ असे समजा.
  $\Gamma$ हा संपूर्ण $\Sigma$-सुसंगत
  असल्याने तो निगामीदृष्ट्या बंद आहे;
  म्हणून $\Gamma \Proves/[\Sigma] \Box \varphi$.
  \ref{nml:com:mod:lem:box2} नुसार
  $\Box^{-1}\Gamma \Proves/[\Sigma] \varphi$.
  \ref{nml:prf:con:prop:consistencyfacts}\ref{nml:prf:con:prop:consistencyfacts-b}
  नुसार $\Box^{-1}\Gamma \cup \{ \lnot\varphi \}$
  हा $\Sigma$-सुसंगत आहे. लिंडेनबाउमच्या
  पूर्वप्रमेयानुसार $\Box^{-1}\Gamma \cup
  \{ \lnot\varphi \} \subseteq \Delta$ असा
  संपूर्ण $\Sigma$-सुसंगत संच~$\Delta$
  असतो. सुसंगततेमुळे $\varphi \notin \Delta$.
\end{proof}


\begin{lem}\label{nml:com:mod:lem:box-iff-diamond}
  $\Gamma$ आणि $\Delta$ हे संपूर्ण
  $\Sigma$-सुसंगत असोत. मग
  $\Box^{-1}\Gamma \subseteq \Delta$
  असेल तेव्हा आणि तेव्हाच
  $\Diamond\Delta \subseteq \Gamma$.
\end{lem}

\begin{proof}
  डावीकडून उजवीकडे जाणारी दिशा: $\Box^{-1}\Gamma
  \subseteq \Delta$ असे समजा आणि
  $\Diamond\varphi \in \Diamond\Delta$
  (म्हणजे $\varphi \in \Delta$) असे धरा.
  $\Diamond\varphi \in \Gamma$ दाखवण्यासाठी
  $\Box\lnot\varphi \notin \Gamma$ दाखवणे
  पुरेसे आहे; कारण मग महत्तमतेमुळे
  $\lnot\Box\lnot \varphi \in \Gamma$.
  आता $\Box\lnot\varphi \in \Gamma$
  असेल, तर गृहीतकानुसार
  $\lnot\varphi \in \Delta$; पण
  $\varphi \in \Delta$ असल्याने हे
  $\Delta$ च्या सुसंगततेच्या
  विरुद्ध आहे. म्हणून आवश्यकतेनुसार
  $\Box\lnot\varphi \notin \Gamma$.

  उजवीकडून डावीकडे जाणारी दिशा: $\Diamond\Delta
  \subseteq \Gamma$ असे समजा.
  प्रतिपरिवर्तन वापरू: $\Box\varphi
  \notin \Gamma$ दाखवण्यासाठी
  $\varphi \notin \Delta$ असे समजा.
  $\varphi \notin \Delta$ असेल,
  तर महत्तमतेमुळे
  $\lnot\varphi \in \Delta$;
  म्हणून गृहीतकानुसार
  $\Diamond\lnot\varphi \in \Gamma$.
  पण सामान्य मोडल तर्कशास्त्रात
  $\Diamond\lnot\varphi$ हे
  $\lnot\Box \varphi$ शी समतुल्य आहे;
  आणि दुसरे सूत्र~$\Gamma$ मध्ये
  असेल, तर सुसंगततेमुळे
  $\Box\varphi \notin\Gamma$,
  हेच दाखवायचे होते.
\end{proof}






\begin{prop}\label{nml:com:mod:prop:diamond}
  $\Gamma$ हा संपूर्ण $\Sigma$-सुसंगत
  असेल, तर $\Diamond \varphi \in \Gamma$
  असेल तेव्हा आणि तेव्हाच
  $\Diamond\Delta \subseteq
  \Gamma$ असा एखादा संपूर्ण
  $\Sigma$-सुसंगत $\Delta$
  असेल की $\varphi \in \Delta$.
\end{prop}

\begin{proof}
$\Gamma$ हा संपूर्ण $\Sigma$-सुसंगत
आहे असे समजा. \Dual{} आणि निगामी बंदतेनुसार
$\Diamond\varphi \in \Gamma$ असेल तेव्हा
आणि तेव्हाच $\lnot\Box\lnot \varphi \in
\Gamma$. $\Gamma$ संपूर्ण
$\Sigma$-सुसंगत असल्याने
\ref{nml:com:ccs:prop:ccs-properties}\ref{nml:com:ccs:prop:ccs-lnot}
नुसार $\lnot\Box\lnot \varphi \in \Gamma$
असेल तेव्हा आणि तेव्हाच
$\Box\lnot \varphi \notin \Gamma$.
\ref{nml:com:mod:prop:box} नुसार
$\Box\lnot \varphi \notin \Gamma$
असेल तेव्हा आणि तेव्हाच
$\Box^{-1}\Gamma \subseteq \Delta$
आणि $\lnot\varphi \notin \Delta$
असा एखादा संपूर्ण $\Sigma$-सुसंगत
$\Delta$ असतो. असा कोणताही~$\Delta$
घ्या. \ref{nml:com:mod:lem:box-iff-diamond}
नुसार $\Box^{-1}\Gamma \subseteq \Delta$
असेल तेव्हा आणि तेव्हाच
$\Diamond\Delta \subseteq \Gamma$.
तसेच
\ref{nml:com:ccs:prop:ccs-properties}\ref{nml:com:ccs:prop:ccs-lnot}
नुसार $\lnot \varphi \notin \Delta$
असेल तेव्हा आणि तेव्हाच
$\varphi \in \Delta$. म्हणून
$\Diamond\varphi \in \Gamma$
असेल तेव्हा आणि तेव्हाच
$\Diamond\Delta \subseteq \Gamma$
आणि $\varphi \in \Delta$ असा एखादा
संपूर्ण $\Sigma$-सुसंगत $\Delta$
असतो.
\end{proof}



\begin{prob}
  $\Gamma$ हा संपूर्ण $\Sigma$-सुसंगत
  असेल, तर पुढील विधान दाखवा:
  $\Diamond \varphi \in \Gamma$
    असेल तेव्हा आणि तेव्हाच
    $\Box^{-1}\Gamma \subseteq \Delta$
    आणि $\varphi \in \Delta$ असा संपूर्ण
    $\Sigma$-सुसंगत $\Delta$ असतो.
  \emph{हे
    \ref{nml:com:mod:lem:box-iff-diamond}
    न वापरता करा}.
\end{prob}













\section{कॅनॉनिकल प्रतिमाने}\label{nml:com:cmd:sec}

सुसंगत मोडल प्रणाली~$\Sigma$ साठीचे \emph{कॅनॉनिकल प्रतिमान}
हे विशिष्ट प्रतिमान~$\mModel{M^\Sigma}$ आहे; त्याची
जगे म्हणजे सर्व संपूर्ण $\Sigma$-सुसंगत संच.
त्याचा प्राप्यता संबंध~$R^\Sigma$ आणि
मूल्यांकन~$V^\Sigma$ यांची व्याख्या अशी केली
जाते की जग~$\Delta$ येथे सत्य असलेली
सूत्रे म्हणजे नेमकी~$\Delta$ मध्ये
असलेली सूत्रे असतात.

\begin{defn}
  $\Sigma$ हे सुसंगत सामान्य मोडल तर्कशास्त्र असू द्या,
  म्हणजे $\Sigma \Proves/ \lfalse$.
  $\Sigma$ साठीचे \emph{कॅनॉनिकल प्रतिमान}
  म्हणजे $\mModel{M}^\Sigma = \tuple{W^\Sigma, R^\Sigma,
    V^\Sigma}$, जेथे:
  \begin{enumerate}
  \item $W^\Sigma = \Setabs{ \Delta }{\Delta \text{
      हा संपूर्ण $\Sigma$-सुसंगत आहे} }$.
  \item $R^\Sigma \Delta\Delta'$ असेल तेव्हा आणि
    तेव्हाच $\Box^{-1}\Delta \subseteq
      \Delta'$.
  \item $V^\Sigma(p) = \Setabs{\Delta \in W^\Sigma}{p \in \Delta}$.
  \end{enumerate}
  सुसंगततेमुळे रिकामा संच $\Sigma$-सुसंगत आहे.
  \ref{nml:com:lin:thm:lindenbaum} नुसार त्याचा संपूर्ण
  $\Sigma$-सुसंगत विस्तार मिळतो; म्हणून $W^\Sigma$ रिकामा नाही.
\end{defn}














\section{सत्यता पूर्वप्रमेय}\label{nml:com:tru:sec}

कॅनॉनिकल प्रतिमान~$\mModel{M^\Sigma}$ ची व्याख्या
अशी केली आहे की $\mSat{M^\Sigma}{\varphi}[\Delta]$
असेल तेव्हा आणि तेव्हाच $\varphi \in \Delta$.
विधानीय चले साठी
$V^\Sigma$ ची व्याख्याच हे थेट दाखवते.
मात्र ही समतुल्यता सर्व सूत्रे साठी
खरी आहे याची पडताळणी करावी लागेल.
हे आपण विगमनाने करू. विगमन-पायरीत
विधानीय कारक असलेल्या सूत्रे साठी
समतुल्यता सिद्ध करावी लागते (येथे
\ref{nml:com:ccs:prop:ccs-properties} वापरतो),
तसेच मोडल कारक असलेल्या सूत्रांसाठीही
(येथे \ref{nml:com:mod:sec} मधील निष्कर्ष वापरतो).

\begin{prop}[सत्यता पूर्वप्रमेय]
  \label{nml:com:tru:prop:truthlemma}
  $\Sigma$ हे सुसंगत सामान्य मोडल तर्कशास्त्र असू द्या.
  प्रत्येक $\Delta \in W^\Sigma$ आणि प्रत्येक सूत्र~$\varphi$ साठी
  $\mSat{M^\Sigma}{\varphi}[\Delta]$ असेल तेव्हा
  आणि तेव्हाच $\varphi \in \Delta$.
\end{prop}

\begin{proof}
  $\varphi$ च्या रचनेवर विगमन करू.
  \begin{enumerate}
    \item \indcase{\varphi}{\lfalse}{$\mSat/{M^\Sigma}{\lfalse}[\Delta]$
        हे \ref{nml:syn:trw:defn:mmodels} नुसार, आणि
        $\lfalse \notin \Delta$ हे
        \ref{nml:com:ccs:prop:ccs-properties}\ref{nml:com:ccs:prop:ccs-lfalse}
        नुसार.}

    \item \indcase{\varphi}{\ltrue}{$\mSat{M^\Sigma}{\ltrue}[\Delta]$
        हे \ref{nml:syn:trw:defn:mmodels} नुसार, आणि
        $\ltrue \in \Delta$ हे
        \ref{nml:com:ccs:prop:ccs-properties}\ref{nml:com:ccs:prop:ccs-ltrue}
        नुसार.}

    \item \indcase{\varphi}{p}{$\mSat{M^\Sigma}{p}[\Delta]$
      असेल तेव्हा आणि तेव्हाच $\Delta \in
      V^\Sigma(p)$, हे \ref{nml:syn:trw:defn:mmodels}
      नुसार. तसेच $\Delta
      \in V^\Sigma(p)$ असेल तेव्हा आणि तेव्हाच
      $p \in \Delta$, हे~$V^\Sigma$ च्या
      व्याख्येनुसार.}

    \item \indcase{\varphi}{\lnot \psi}{
        $\mSat{M^\Sigma}{\lnot \psi}[\Delta]$
          असेल तेव्हा आणि तेव्हाच
          $\mSat/{M^\Sigma}{\psi}[\Delta]$
          (\ref{nml:syn:trw:defn:mmodels} नुसार);
          असेल तेव्हा आणि तेव्हाच $\psi \notin \Delta$
          (विगमन गृहीतकानुसार);
          असेल तेव्हा आणि तेव्हाच
          $\lnot \psi \in \Delta$
          (\ref{nml:com:ccs:prop:ccs-properties}\ref{nml:com:ccs:prop:ccs-lnot}
          नुसार).}

    \item \indcase{\varphi}{\psi \land \chi}{
        $\mSat{M^\Sigma}{\psi \land
            \chi}[\Delta]$ असेल तेव्हा आणि तेव्हाच
          $\mSat{M^\Sigma}{\psi}[\Delta]$ आणि
          $\mSat{M^\Sigma}{\chi}[\Delta]$
          (\ref{nml:syn:trw:defn:mmodels} नुसार);
          असेल तेव्हा आणि तेव्हाच $\psi \in \Delta$
          आणि $\chi \in \Delta$
          (विगमन गृहीतकानुसार);
          असेल तेव्हा आणि तेव्हाच $\psi \land \chi \in
          \Delta$
          (\ref{nml:com:ccs:prop:ccs-properties}\ref{nml:com:ccs:prop:ccs-land}
          नुसार).}

    \item \indcase{\varphi}{\psi \lor \chi}{
        $\mSat{M^\Sigma}{\psi \lor
            \chi}[\Delta]$ असेल तेव्हा आणि तेव्हाच
          $\mSat{M^\Sigma}{\psi}[\Delta]$ किंवा
          $\mSat{M^\Sigma}{\chi}[\Delta]$
          (\ref{nml:syn:trw:defn:mmodels} नुसार);
          असेल तेव्हा आणि तेव्हाच $\psi \in \Delta$
          किंवा $\chi \in \Delta$
          (विगमन गृहीतकानुसार);
          असेल तेव्हा आणि तेव्हाच $\psi \lor \chi \in
          \Delta$
          (\ref{nml:com:ccs:prop:ccs-properties}\ref{nml:com:ccs:prop:ccs-lor}
          नुसार).}

    \item \indcase{\varphi}{\psi \lif \chi}{
        $\mSat{M^\Sigma}{\psi \lif
            \chi}[\Delta]$ असेल तेव्हा आणि तेव्हाच
          $\mSat/{M^\Sigma}{\psi}[\Delta]$ किंवा
          $\mSat{M^\Sigma}{\chi}[\Delta]$
          (\ref{nml:syn:trw:defn:mmodels} नुसार);
          असेल तेव्हा आणि तेव्हाच $\psi \notin \Delta$
          किंवा $\chi \in \Delta$
          (विगमन गृहीतकानुसार);
          असेल तेव्हा आणि तेव्हाच $\psi \lif \chi
          \in \Delta$
          (\ref{nml:com:ccs:prop:ccs-properties}\ref{nml:com:ccs:prop:ccs-lif}
          नुसार).}

    \item \indcase{\varphi}{\psi \liff \chi}{
        $\mSat{M^\Sigma}{\psi \liff
            \chi}[\Delta]$ असेल तेव्हा आणि तेव्हाच
          $\mSat{M^\Sigma}{\psi}[\Delta]$ आणि
          $\mSat{M^\Sigma}{\chi}[\Delta]$
          ही दोन्ही, किंवा
          $\mSat/{M^\Sigma}{\psi}[\Delta]$ आणि
          $\mSat/{M^\Sigma}{\chi}[\Delta]$
          ही दोन्ही
          (\ref{nml:syn:trw:defn:mmodels} नुसार);
          असेल तेव्हा आणि तेव्हाच $\psi \in \Delta$
          आणि $\chi \in \Delta$ ही दोन्ही, किंवा
          $\psi \notin \Delta$ आणि $\chi \notin
          \Delta$ ही दोन्ही
          (विगमन गृहीतकानुसार);
          असेल तेव्हा आणि तेव्हाच $\psi \liff \chi \in
          \Delta$
          (\ref{nml:com:ccs:prop:ccs-properties}\ref{nml:com:ccs:prop:ccs-liff}
          नुसार).}

    \item \indcase{\varphi}{\Box \psi}{
        प्रथम
          $\mSat{M^\Sigma}{\Box \psi}[\Delta]$
          असे समजा.
          \ref{nml:syn:trw:defn:mmodels} नुसार
          $R^\Sigma \Delta\Delta'$ असलेल्या
          प्रत्येक $\Delta'$ साठी
          $\mSat{M^\Sigma}{\psi}[\Delta']$.
          विगमन गृहीतकानुसार
          $R^\Sigma \Delta\Delta'$ असलेल्या
          प्रत्येक $\Delta'$ साठी $\psi \in
          \Delta'$. $R^\Sigma$ च्या
          व्याख्येनुसार $\Box^{-1} \Delta
          \subseteq \Delta'$ असलेल्या प्रत्येक
          $\Delta'$ साठी $\psi \in \Delta'$.
          \ref{nml:com:mod:prop:box} नुसार
          $\Box \psi \in \Delta$.

          आता $\Box \psi \in \Delta$ असे समजा.
          $R^\Sigma \Delta\Delta'$ म्हणजेच
          $\Box^{-1}\Delta \subseteq \Delta'$
          असलेला कोणताही $\Delta' \in W^\Sigma$
          घ्या. $\Box \psi \in \Delta$
          असल्याने $\psi \in \Box^{-1} \Delta$.
          परिणामी $\psi \in \Delta'$.
          विगमन गृहीतकानुसार
          $\mSat{M^\Sigma}{\psi}[\Delta']$.
          $R^\Sigma \Delta\Delta'$ असलेला
          $\Delta'$ स्वेच्छ असल्याने
          $R^\Sigma \Delta\Delta'$ असलेल्या
          सर्व $\Delta' \in W^\Sigma$ साठी
          $\mSat{M^\Sigma}{\psi}[\Delta']$.
          \ref{nml:syn:trw:defn:mmodels} नुसार
          $\mSat{M^\Sigma}{\Box
            \psi}[\Delta]$.}

    \item \indcase{\varphi}{\Diamond \psi}{
        प्रथम
          $\mSat{M^\Sigma}{\Diamond \psi}[\Delta]$
          असे समजा.
          \ref{nml:syn:trw:defn:mmodels} नुसार
          $R^\Sigma \Delta\Delta'$ असलेल्या
          एखाद्या $\Delta'$ साठी
          $\mSat{M^\Sigma}{\psi}[\Delta']$.
          विगमन गृहीतकानुसार
          $R^\Sigma \Delta\Delta'$ असलेल्या
          एखाद्या $\Delta'$ साठी
          $\psi \in \Delta'$.
          $R^\Sigma$ च्या व्याख्येनुसार,
          $\Box^{-1} \Delta
            \subseteq \Delta'$ असलेल्या
            एखाद्या $\Delta'$ साठी
            $\psi \in \Delta'$.
            \ref{nml:com:mod:lem:box-iff-diamond} नुसार,
          $\Diamond \Delta' \subseteq \Delta$
          असलेल्या एखाद्या $\Delta'$ साठी
          $\psi \in \Delta'$.
          $\psi \in \Delta'$ असल्याने
          $\Diamond \psi \in \Diamond
          \Delta'$; म्हणून
          $\Diamond \psi \in \Delta$.

          आता $\Diamond \psi \in \Delta$
          असे समजा. \ref{nml:com:mod:prop:diamond}
          नुसार $\Diamond \Delta' \subseteq
          \Delta$ आणि $\psi \in \Delta'$
          असलेला संपूर्ण $\Sigma$-सुसंगत
          $\Delta' \in W^\Sigma$ असतो.
          \ref{nml:com:mod:lem:box-iff-diamond}
            नुसार $\Box^{-1}
            \Delta \subseteq \Delta'$
            आणि $\psi \in \Delta'$ असलेला
            $\Delta' \in W^\Sigma$ असतो.
          $R^\Sigma$ च्या व्याख्येनुसार
          $R^\Sigma \Delta\Delta'$;
          म्हणून $R^\Sigma
          \Delta\Delta'$ आणि
          $\psi \in \Delta'$ असलेला
          $\Delta' \in W^\Sigma$ असतो.
          विगमन गृहीतकानुसार $\mSat{M^\Sigma}{\psi}[\Delta']$.
          \ref{nml:syn:trw:defn:mmodels}
          नुसार $\mSat{M^\Sigma}{\Diamond
            \psi}[\Delta]$.}
  \end{enumerate}
\end{proof}















\section{\Log{K} चे निर्धारण आणि संपूर्णता}\label{nml:com:cmk:sec}

आता संपूर्णता सिद्ध करण्यासाठी कॅनॉनिकल प्रतिमान
वापरता येईल. $\Sigma$ साठीच्या कॅनॉनिकल
प्रतिमानात सत्य असलेली सूत्रे म्हणजे
नेमकी $\Sigma$-निष्पन्न करता येण्याजोगे सूत्रे असतात;
या तथ्यावरून संपूर्णता मिळते. हा गुणधर्म
असलेले प्रतिमान~$\Sigma$ चे
\emph{निर्धारण करते} असे म्हणतात.

\begin{defn}
  सर्व सूत्रांपैकी प्रत्येक~$\varphi$ साठी
  $\mSat{M}{\varphi}$ असेल तेव्हा आणि तेव्हाच
  $\Sigma \Proves \varphi$ असेल, तर
  प्रतिमान~$\mModel{M}$ हे सामान्य मोडल
  तर्कशास्त्र~$\Sigma$ चे
  \emph{निर्धारण करते}.
\end{defn}

\begin{thm}[निर्धारण]\label{nml:com:cmk:thm:determination}
   $\Sigma$ हे सुसंगत सामान्य मोडल तर्कशास्त्र असू द्या. मग
   $\mSat{M^\Sigma}{\varphi}$ असेल तेव्हा
   आणि तेव्हाच $\Sigma \Proves \varphi$.
\end{thm}

\begin{proof}
  $\mSat{M^\Sigma}{\varphi}$ असेल, तर प्रत्येक
  संपूर्ण $\Sigma$-सुसंगत $\Delta$ साठी
  $\mSat{M^\Sigma}{\varphi}[\Delta]$.
  म्हणून सत्यता पूर्वप्रमेयानुसार प्रत्येक
  संपूर्ण $\Sigma$-सुसंगत $\Delta$ साठी
  $\varphi \in \Delta$. त्यामुळे
  \ref{nml:com:lin:cor:provability-characterization}
  वरून ($\Gamma = \emptyset$ घेऊन)
  $\Sigma \Proves \varphi$.

  उलट $\Sigma \Proves \varphi$ असेल, तर
  \ref{nml:com:ccs:prop:ccs-properties}\ref{nml:com:ccs:prop:ccs-closed}
  नुसार प्रत्येक संपूर्ण $\Sigma$-सुसंगत
  $\Delta$ मध्ये~$\varphi$ असते. म्हणून
  सत्यता पूर्वप्रमेयानुसार प्रत्येक~$\Delta \in
  W^\Sigma$ साठी
  $\mSat{M^\Sigma}{\varphi}[\Delta]$;
  म्हणजे $\mSat{M^\Sigma}{\varphi}$.
\end{proof}

निर्दोषतेनुसार \Log{K} सुसंगत आहे, कारण असत्य सूत्र वैध नाही.
\Log{K} साठीचे कॅनॉनिकल प्रतिमान~\Log{K}
चे निर्धारण करत असल्याने
\Log{K} ची संपूर्णता लगेच निष्पन्नतेने मिळते:

\begin{cor}\label{nml:com:cmk:cor:Kcomplete}
  मूलभूत मोडल तर्कशास्त्र~\Log{K}
सर्व प्रतिमानांच्या वर्गाच्या सापेक्ष
संपूर्ण आहे; म्हणजे
$\Entails \varphi$ असेल, तर $\Log{K}
\Proves \varphi$.
\end{cor}

\begin{proof}
  प्रतिपरिवर्तनाने: $\Log{K} \Proves/ \varphi$
  असेल, तर निर्धारणावरून
  $\mSat/{M^{\Log{K}}}{\varphi}$;
  म्हणून~$\varphi$ वैध नाही.
\end{proof}

सर्वसाधारणपणे एखादी प्रणाली~$\Sigma$
प्रतिमानांच्या एखाद्या वर्गाच्या सापेक्ष
संपूर्ण आहे हे सिद्ध करण्यासाठी
केवळ निर्धारण पुरेसे नसते. उदाहरणार्थ,
\Log{KTB4} ची परावर्ती, सममित आणि
संक्रमक प्रतिमानांच्या वर्गाच्या सापेक्ष
संपूर्णता विचारात घ्या. प्रणाली~$\Sigma$
साठीच्या कॅनॉनिकल प्रतिमानातून ही संपूर्णता सिद्ध करायची असेल,
तर ते प्रतिमान त्या वर्गातील सदस्य आहे हेही दाखवावे लागते.
कॅनॉनिकल प्रतिमानाच्या रचनेवरून हे
स्पष्टपणे मिळत नाही---काही वेळा
तर ते खरेही नसते!













\section{चौकट-संपूर्णता}\label{nml:com:fra:sec}

दिलेल्या तर्कशास्त्राच्या कॅनॉनिकल प्रतिमानाला
संबंधित चौकट-गुणधर्म आहे हे दाखवल्यावर
\Log{K} साठीचे संपूर्णता प्रमेय इतर
मोडल प्रणालींनाही लागू करता येते.

\begin{thm}\label{nml:com:fra:thm:completeframeprops}
  सुसंगत सामान्य मोडल तर्कशास्त्र~$\Sigma$ मध्ये
  \ref{nml:com:fra:tab:correspondencetable} च्या
  डाव्या बाजूच्या सूत्रांपैकी एखादे असेल,
  तर~$\Sigma$ साठीच्या कॅनॉनिकल प्रतिमानाला
  उजव्या बाजूचा संबंधित गुणधर्म असतो.
\end{thm}

\begin{table}[htp]
  \centering
    \begin{tabular}{| l || l |}
      \hline
      \emph{$\Sigma$ मध्ये \dots असेल, तर} & \emph{$\Sigma$ साठीचे
        कॅनॉनिकल प्रतिमान \dots असते:} \\
      \hline \hline
      \Ax{D}:\;\; $\Box\varphi \lif \Diamond \varphi$ & \quad सर्वत्र उत्तरवर्ती; \\
      \hline
      \Ax{T}:\;\; $\Box\varphi \lif \varphi$ &\quad  परावर्ती;\\
      \hline
      \Ax{B}: \;\;$\varphi \lif \Box\Diamond\varphi$ &\quad  सममित; \\
      \hline
      \Ax{4}: \;\; $\Box\varphi \lif \Box\Box\varphi$ & \quad संक्रमक; \\
      \hline
      \Ax{5}: \;\; $\Diamond \varphi \lif \Box\Diamond\varphi$& \quad युक्लिडी.\\
      \hline
    \end{tabular}
    \caption{मूलभूत अनुरूपता तथ्ये.}
    \label{nml:com:fra:tab:correspondencetable}
\end{table}

\begin{proof}
  यांपैकी प्रत्येक प्रकरण स्वतंत्रपणे पाहू.

  $\Sigma$ मध्ये \Ax{D} आहे असे समजा आणि
  $\Delta \in W^\Sigma$ घ्या.
  $R^\Sigma \Delta\Delta'$ असा
  $\Delta'$ अस्तित्वात आहे हे दाखवायचे आहे.
  त्यासाठी $\Box^{-1}\Delta$ हा
  $\Sigma$-सुसंगत आहे हे दाखवणे पुरेसे;
  कारण मग लिंडेनबाउमच्या पूर्वप्रमेयानुसार
  त्याचा विस्तार असलेला संपूर्ण
  $\Sigma$-सुसंगत संच
  $\Delta' \supseteq \Box^{-1}\Delta$
  असतो आणि $R^\Sigma$ च्या व्याख्येनुसार

  $R^\Sigma \Delta\Delta'$.
  म्हणून व्याघातासाठी $\Box^{-1}\Delta$
  हा $\Sigma$-सुसंगत \emph{नाही}
  असे समजा; म्हणजे
  $\Box^{-1}\Delta \Proves[\Sigma] \lfalse$.
  \ref{nml:com:mod:lem:box2} नुसार
  $\Delta \Proves[\Sigma] \Box\lfalse$;
  आणि $\Sigma$ मध्ये \Ax{D} असल्याने
  $\Delta \Proves[\Sigma] \Diamond\lfalse$
  हेही मिळते. पण $\Sigma$ सामान्य असल्याने
  $\Sigma \Proves \lnot\Diamond \lfalse$
  (\ref{nml:prf:nor:prop:notDiamondBot});
  म्हणून $\Delta
  \Proves[\Sigma] \lnot\Diamond \lfalse$
  हेही मिळते. हा~$\Delta$ च्या
  सुसंगततेशी व्याघात आहे.

  आता $\Sigma$ मध्ये \Ax{T} आहे
  असे समजा आणि $\Delta \in W^\Sigma$
  घ्या. $R^\Sigma \Delta\Delta$
  दाखवायचे आहे; म्हणजेच

  $\Box^{-1}\Delta \subseteq \Delta$.
  पण $\Box\varphi \in \Delta$ असेल,
  तर \Ax{T} नुसार $\varphi \in \Delta$;
  हेच अपेक्षित होते.

  आता $\Sigma$ मध्ये \Ax{B} आहे
  असे समजा आणि $\Delta$, $\Delta' \in
  W^\Sigma$ साठी $R^\Sigma
  \Delta\Delta'$ असे धरा.
  $R^\Sigma \Delta'\Delta$
  दाखवायचे आहे; म्हणजे
  $\Box^{-1}\Delta'
    \subseteq \Delta$.
    \ref{nml:com:mod:lem:box-iff-diamond} नुसार
    ही अट पुढील अटीशी समतुल्य आहे:
  $\Diamond\Delta \subseteq \Delta'$.
  म्हणून $\varphi \in \Delta$ असे
  समजा. \Ax{B} नुसार
  $\Box\Diamond\varphi \in \Delta$.
  $R^\Sigma \Delta\Delta'$ या
  गृहीतकावरून
  $\Box^{-1}\Delta \subseteq \Delta'$;
  म्हणून आवश्यकतेनुसार
  $\Diamond\varphi \in \Delta'$.

  आता $\Sigma$ मध्ये \Ax{4} आहे असे
  समजा आणि $R^\Sigma \Delta_1\Delta_2$
  व $R^\Sigma \Delta_2\Delta_3$
  असे धरा. $R^\Sigma
  \Delta_1\Delta_3$ दाखवायचे आहे.
  गृहीतकावरून
  $\Box^{-1}\Delta_1 \subseteq \Delta_2$
  आणि $\Box^{-1}\Delta_2 \subseteq \Delta_3$
  या दोन्ही अटी मिळतात.
  $R^\Sigma \Delta_1\Delta_3$
  साठी $\Box^{-1}\Delta_1
  \subseteq \Delta_3$ दाखवणे पुरेसे.
  म्हणून $\psi \in \Box^{-1}\Delta_1$
  घ्या; म्हणजे $\Box\psi \in \Delta_1$.
  \Ax{4} नुसार $\Box\Box\psi \in
  \Delta_1$ हेही मिळते. गृहीतकांवरून
  आधी $\Box\psi \in \Delta_2$
  आणि नंतर $\psi \in \Delta_3$,
  हेच अपेक्षित होते.

  आता $\Sigma$ मध्ये \Ax{5} आहे
  असे समजा आणि $R^\Sigma
  \Delta_1\Delta_2$ व
  $R^\Sigma \Delta_1\Delta_3$
  असे धरा. $R^\Sigma
  \Delta_2\Delta_3$ दाखवायचे आहे.
  पहिल्या गृहीतकावरून
  $\Box^{-1}\Delta_1 \subseteq \Delta_2$

  मिळते; दुसरे गृहीतक
  $\Diamond\Delta_3 \subseteq \Delta_1$
  याच्याशी समतुल्य आहे
  \ref{nml:com:mod:lem:box-iff-diamond} नुसार.
  $R^\Sigma \Delta_2\Delta_3$
  दाखवण्यासाठी
    \ref{nml:com:mod:lem:box-iff-diamond} नुसार,
    पुढील अट दाखवणे पुरेसे:
  $\Diamond\Delta_3 \subseteq \Delta_2$.
  म्हणून $\Diamond\varphi \in
  \Diamond\Delta_3$ घ्या;
  म्हणजे $\varphi \in \Delta_3$.
  दुसऱ्या गृहीतकानुसार
  $\Diamond\varphi \in \Delta_1$;
  \Ax{5} नुसार
  $\Box\Diamond\varphi \in \Delta_1$
  हेही मिळते. आता पहिल्या
  गृहीतकावरून $\Diamond\varphi \in
  \Delta_2$, हेच अपेक्षित होते.
  \end{proof}

या निष्कर्षांवरून अनेक प्रणालींसाठी
संपूर्णतेची निष्पन्नता मिळते.
उदाहरणार्थ, $\Log{S5} = \Log{KT5} =
\Log{KTB4}$ हे सर्व परावर्ती युक्लिडी
प्रतिमानांच्या वर्गाच्या सापेक्ष संपूर्ण
आहे, हे आपल्याला माहीत आहे.
हा वर्ग सर्व परावर्ती, सममित आणि
संक्रमक प्रतिमानांच्या वर्गाइतकाच आहे.

\begin{thm}\label{nml:com:fra:thm:generaldet}
  $\mClass{C}_\Ax{D}$,
  $\mClass{C}_\Ax{T}$,
  $\mClass{C}_\Ax{B}$,
  $\mClass{C}_\Ax{4}$ आणि
  $\mClass{C}_\Ax{5}$
  हे अनुक्रमे सर्व सर्वत्र उत्तरवर्ती,
  परावर्ती, सममित, संक्रमक आणि
  युक्लिडी प्रतिमानांचे वर्ग असू द्या.
  मग \Ax{D}, \Ax{T}, \Ax{B},
  \Ax{4} आणि \Ax{5} यांमधील
  कोणत्याही योजना $\varphi_1$, \dots, $\varphi_n$
  साठी प्रणाली
  $\Log{K}\varphi_1 \dots \varphi_n$
  हिचे निर्धारण प्रतिमानांचा वर्ग
  $\mClass{C} = \mClass{C}_{\varphi_1} \cap \dots
  \cap \mClass{C}_{\varphi_n}$ करतो.
\end{thm}

\begin{prop}
  $\Sigma$ हे सुसंगत सामान्य मोडल तर्कशास्त्र
  असू द्या. मग:
  \begin{enumerate}
  \item\label{nml:com:fra:prop:anotherfive-a}
    $\Sigma$ मध्ये $\Diamond\varphi \lif \Box
    \varphi$ ही योजना असेल, तर~$\Sigma$
    साठीचे कॅनॉनिकल प्रतिमान
    आंशिक फलनीय असते.
  \item $\Sigma$ मध्ये
    $\Diamond\varphi \liff \Box
    \varphi$ ही योजना असेल, तर~$\Sigma$
    साठीचे कॅनॉनिकल प्रतिमान
    फलनीय असते.
  \item $\Sigma$ मध्ये
    $\Box\Box\varphi \lif \Box
    \varphi$ ही योजना असेल, तर~$\Sigma$
    साठीचे कॅनॉनिकल प्रतिमान
    दुर्बल सघन असते.
  \end{enumerate}
  (या चौकट-गुणधर्मांच्या व्याख्यांसाठी
  \ref{nml:frd:acc:tab:anotherfive} पाहा.)
\end{prop}

\begin{proof}
  \begin{enumerate}
  \item $\Sigma$ मध्ये
    $\Diamond \varphi \lif \Box \varphi$
    ही योजना आहे असे समजा.
    $R^\Sigma$ आंशिक फलनीय आहे हे
    दाखवण्यासाठी कोणत्याही
    $\Delta_1$, $\Delta_2$,
    $\Delta_3 \in W^\Sigma$
    साठी $R^\Sigma \Delta_1\Delta_2$
    आणि $R^\Sigma \Delta_1\Delta_3$
    असतील, तर $\Delta_2=\Delta_3$
    हे सिद्ध करायचे आहे.
    $R^\Sigma \Delta_1\Delta_2$
    असल्याने $\Box^{-1}\Delta_1
    \subseteq \Delta_2$;
    आणि $R^\Sigma \Delta_1\Delta_3$
    असल्याने $\Box^{-1}\Delta_1
    \subseteq \Delta_3$
    .
    $\Delta_2=\Delta_3$ साठी
    $\Delta_2 \subseteq \Delta_3$
    आणि $\Delta_3 \subseteq \Delta_2$
    हे दोन्ही समावेश सिद्ध करणे पुरेसे.
    पहिल्या समावेशासाठी
    $\varphi \in \Delta_2$ घ्या;
    मग $\Diamond\varphi \in \Delta_1$.
    या योजनेवरून आणि~$\Delta_1$ च्या
    निगामी बंदतेवरून
    $\Box\varphi \in \Delta_1$ हेही मिळते.
    $R^\Sigma \Delta_1\Delta_3$
    या गृहीतकावरून $\varphi \in \Delta_3$.
    दुसरा समावेशही तसाच सिद्ध होतो.

  \item $\Diamond\varphi \liff \Box\varphi$ या योजनेतून
    $\Box\varphi \lif \Diamond\varphi$, म्हणजे \Ax{D}, मिळते.
    म्हणून हे \ref{nml:com:fra:prop:anotherfive-a}
    आणि \ref{nml:com:fra:thm:completeframeprops}
    मधील सर्वत्र उत्तरवर्तित्वाच्या
    सिद्धतेवरून लगेच मिळते.

  \item $\Sigma$ मध्ये
    $\Box\Box\varphi \lif \Box\varphi$
    ही योजना आहे असे समजा.
    $R^\Sigma$ दुर्बल सघन आहे हे
    दाखवण्यासाठी $R^\Sigma
    \Delta_1\Delta_2$ घ्या.
    $R^\Sigma \Delta_1\Delta_3$
    आणि $R^\Sigma \Delta_3\Delta_2$
    असा संपूर्ण $\Sigma$-सुसंगत
    संच~$\Delta_3$ अस्तित्वात आहे
    हे दाखवायचे आहे. असे ठेवू:
    \[    \Gamma = \Box^{-1}\Delta_1 \cup \Diamond\Delta_2.
    \]
    $\Gamma$ हा $\Sigma$-सुसंगत
    आहे हे दाखवणे पुरेसे. मग
    लिंडेनबाउमच्या पूर्वप्रमेयानुसार
    त्याचा विस्तार असा संपूर्ण
    $\Sigma$-सुसंगत संच~$\Delta_3$
    मिळतो की $\Box^{-1}\Delta_1
    \subseteq \Delta_3$ आणि
    $\Diamond\Delta_2 \subseteq \Delta_3$.
    म्हणजे $R^\Sigma \Delta_1\Delta_3$

    आणि $R^\Sigma \Delta_3\Delta_2$
    (
      \ref{nml:com:mod:lem:box-iff-diamond} नुसार).

    व्याघातासाठी $\Gamma$ सुसंगत
    नाही असे समजा. $\Box^{-1}\Delta_1\subseteq\Delta_2$
    आणि $\Delta_2$ सुसंगत असल्याने फक्त पहिल्या संचातील
    गृहीतके वापरून व्याघात मिळू शकत नाही. म्हणून
    पुढील finite witness मध्ये $n\geq0$ आणि $m\geq1$ निवडता येतात.
    मग $\Box\varphi_1$, \dots,
    $\Box\varphi_n \in \Delta_1$
    आणि $\psi_1$, \dots,~$\psi_m \in \Delta_2$
    अशी सूत्रे असतात की
    \[\varphi_1, \dots,
    \varphi_n, \Diamond\psi_1, \dots, \Diamond\psi_m \Proves[\Sigma]
    \lfalse.\]
    $\Diamond (\psi_1 \land \dots \land \psi_m) \to
    (\Diamond\psi_1 \land \dots \land \Diamond\psi_m)$
    हे प्रत्येक सामान्य मोडल तर्कशास्त्रात
    निष्पन्न करता येण्याजोगे असल्याने
    पुढीलप्रमाणे युक्तिवाद करून
    आपण~$\Delta_2$ च्या
    सुसंगततेशी व्याघात मिळवतो:
    \begin{align*}
      \varphi_1, \dots, \varphi_n, & \Diamond\psi_1, \dots,
      \Diamond\psi_m \Proves[\Sigma] \lfalse \\
      \varphi_1, \dots,\varphi_n & \Proves[\Sigma] (\Diamond\psi_1 \land
      \dots \land \Diamond\psi_m) \lif \lfalse\\
      & \qquad\text{निगमन प्रमेयानुसार}\\
      & \qquad\text{\ref{nml:prf:prp:prop:derivabilityfacts}\ref{nml:prf:prp:prop:derivabilityfacts-deduction}, आणि \Taut} \\
      \varphi_1, \dots,\varphi_n & \Proves[\Sigma]
      \Diamond(\psi_1 \land
      \dots \land \psi_m) \lif \lfalse\\
      & \qquad \text{$\Sigma$ सामान्य असल्याने} \\
      \varphi_1, \dots,\varphi_n & \Proves[\Sigma]
      \lnot\Diamond (\psi_1 \land \dots \land \psi_m)\\
      & \qquad\text{\PL नुसार} \\
      \varphi_1, \dots,\varphi_n & \Proves[\Sigma]
      \Box\lnot (\psi_1 \land \dots \land \psi_m)\\
      & \qquad\text{$\lnot\Diamond$ साठी $\Box\lnot$} \\
      \Box\varphi_1, \dots,\Box\varphi_n
      & \Proves[\Sigma] \Box\Box \lnot (\psi_1 \land \dots \land \psi_m)\\
      & \qquad\text{\ref{nml:com:mod:lem:box1} नुसार} \\
      \Box\varphi_1, \dots,\Box\varphi_n
      & \Proves[\Sigma]
      \Box\lnot (\psi_1 \land \dots \land \psi_m)\\
      &\qquad\text{$\Box\Box\varphi \lif \Box\varphi$ योजनेनुसार} \\
      \Delta_1 & \Proves[\Sigma]
      \Box\lnot (\psi_1 \land \dots \land \psi_m)\\
      &\qquad\text{एकस्वनिकतेनुसार, \ref{nml:prf:prp:prop:derivabilityfacts}\ref{nml:prf:prp:prop:derivabilityfacts-monotonicity}} \\
      & \Box\lnot (\psi_1 \land \dots \land \psi_m) \in \Delta_1\\
      &\qquad\text{निगामी बंदतेनुसार}; \\
      & \lnot (\psi_1 \land \dots \land \psi_m) \in \Delta_2\\
      &\qquad \text{कारण }
      R^\Sigma \Delta_1\Delta_2.
    \end{align*}
  \end{enumerate}
\end{proof}

या उदाहरणांवरून प्रत्येक मोडल
प्रणाली~$\Sigma$ \emph{संपूर्ण}
असेल असे वाटू शकते: म्हणजे~$\Sigma$
ची सर्व प्रमेये ज्या प्रत्येक चौकटीत
वैध आहेत, त्या प्रत्येक चौकटीत
वैध असलेले प्रत्येक सूत्रही
ती प्रणाली सिद्ध करत असेल.
पण दुर्दैवाने या अर्थाने संपूर्ण
नसलेल्या अनेक प्रणाली आहेत.



